%But : introduire l'action du groupe de classe sur les isogénie

\section{Isogénies et idéaux d'ordres quadratiques imaginaires}

\subsection{Rappels de théorie des nombres}
%pas/peu de preuves
 Le stage a débuté par une mise à niveau en théorie des nombres, dont on énonce ici les résultats principaux utiles par la suite. On se réfère à \cite{Cox}[Ch 5, 7].

\subsubsection{Ordres quadratiques imaginaires}

Dans cette section, on énonce des résultats généraux de théorie des nombres. On se réfère à \cite{Cox}[Ch 5, 7]. On ne s'intéresse qu'à certaines extensions de degrés $2$ de $\Q$.

\begin{Def}
	
	Un corps $\K$ est appelé corps quadratique imaginaire si $\K = \Q(\sqrt{D})$ avec $D$ entier négatif sans facteur carré. Un ordre de $\K$ est un sous anneaux unitaire de $\K$, finement engendré comme $\Z$-module et contenant une $\Q$-base de $\K$. Un tel anneaux est appelé un ordre quadratique imaginaire. 
\end{Def}

\begin{Prop}\cite{Cox}[Prop 5.3, 7.2]
	
	Soit $\OR$ un ordre de $\K$. Alors $\OR$ est un $\Z$-module libre de rang $\left[ \K : \Q\right] $. De plus $\K = Frac(\OR)$.
	On note $\OR_{\K}$ l'anneau des entiers algébriques d'un corps de nombre $\K$. C'est un ordre de $\K$, maximal pour l'inclusion et contenant tout les ordres de $\K$. 
\end{Prop}

\begin{Def}
	Le discriminant d'un corps quadratique imaginaire $\K = \Q(\sqrt{D})$, noté $d_{\K}$, est défini par :
	\begin{enumerate}
		\item $d_{\K} = D$ si $D = 1 [4]$,
		\item $d_{\K} = 4D$ sinon.
	\end{enumerate}
	En particulier $d_{\K} = 0 ,\;1[4]$. On pose de plus :
	\begin{equation*}
		w_{\K} = \frac{d_{\K} + \sqrt{d_{\K}}}{2}
	\end{equation*}
\end{Def}


\begin{Prop}\cite{Cox}[Prop 5.14, 7.2] 
	
	Soit $\K$ un corps quadratique imaginaire. Alors $\OR_{\K} = \Z\left[ w_{\K}\right] $. De plus si $\K = \Q(\sqrt{D})$ avec $D$ sans facteur carré, alors
	\begin{enumerate}
		\item Si $D \neq 1 [4]$, $\OR_{\K} = \Z\left[ \sqrt{D}\right] $
		\item Sinon, $\OR_{\K} = \Z\left[ \frac{1 + \sqrt{D}}{2}\right] $
	\end{enumerate} 
	Soit $\OR$ un ordre de $\K = \Q(\sqrt{D})$ un corps quadratique imaginaire. $\OR\subset\OR_{\K}$ et on note $f = \left[ \OR_{\K} : \OR\right]$. Alors $\OR = \Z\left[ fw_{\K}\right] $.  
\end{Prop}

\begin{Def}
	Lorsque $\K$ est un corps quadratique imaginaire, si $\OR$ est un ordre de $\K$, $f = \left[ \OR_{\K} : \OR\right]$ est appelé le conducteur de $\OR$, et $\Delta = f^2d_{\K}$ son discriminant. L'ordre de $\K$ de discriminant $\Delta$ est noté $\OR_{\Delta}$
\end{Def}

La notion de conducteur est multiplicative au sens suivant:

\begin{Prop}
	Soit $\OR$, $\OR'$ deux ordres d'un $\K$ quadratique imaginaire. Soit $f$ et $f'$ leurs conducteurs. Si $\OR'\subset\OR$ alors $f$ divise $f'$
\end{Prop}

On peut reformuler ces résultats en terme de réseau de $\C$.

\begin{Def}
	Un réseau $\Lambda$ de $\C$ est un ensemble de la forme $\left\lbrace nz_{1} + mz_{2} , n,m\in\Z\right\rbrace$ où $z_{1},\; z_{2}\in\C$ sont linéairement indépendant sur $\R$. On note alors $\Lambda = \left[ z_{1}, z_{2}\right] $
\end{Def}

\begin{Rem}
	Les ordres quadratiques imaginaires sont donc des réseaux de $\C$, qui contiennent $1$ : $\OR_{\K} = \left[1 , w_{\K}\right]$, et si $\OR$ est l'ordre de conducteur $f$, alors $\OR = \left[ 1, fw_{\K}\right]$.
\end{Rem}

Les ordres de $\K$ ne sont pas des anneaux factoriels en général. On va cependant obtenir des résultats de factorisations au niveau des idéaux. En particulier $\OR_{\K}$ est un anneaux de Dedekind, un type d'anneaux qui admet une unique factorisation au niveau des idéaux par des idéaux premiers.

\begin{Prop}\cite{Cox}[Prop 5.4, Ex 5.1, Ex 7.4]
	
	 Soit $\LR\subset\OR$ un idéal d'un ordre d'un corps quadratique imaginaire $\K$.
	\begin{enumerate}
		\item Il existe $a\in\Z$ appartenant à $\LR$
		\item $\LR$ est un $\Z$-module libre de rang $2$
		\item Le quotient $\OR/\LR$ est fini.
		\item Si $\LR$ est premier, alors $\LR$ est maximal. 
	\end{enumerate}
\end{Prop}

\begin{Def}
	La norme d'un idéal $\LR\subset\OR$ d'un ordre d'un corps quadratique imaginaire $\K$ est $N(\LR) = |\OR/\LR|$.
\end{Def}

\begin{Coro} \textbf{TODO citation}
	
	Soit $\OR$ un ordre de $\K = \Q(\sqrt{D})$ un corps quadratique imaginaire de conducteur $f$. Tout idéal $\LR$ de $\OR$ est un réseaux de $\C$ de la forme $\left[ a, b + cfw_{\K}\right]$ avec $a | N(\LR)$, $b$, $c\in\Z$, $c\neq 0$ et $0 \leq b< a$. 	
\end{Coro}

\begin{proof}
	$\LR\cap\Z$ est un idéal non nul de $\Z$, donc principal engendré par un entier $a > 0$. Alors l'anneau fini $\OR/\LR$ est de caractéristique $a$, et $a | N(\LR)$. Soit $\alpha\in\LR$ linéairement indépendant de $a$. Comme $\LR$ est un sous ensemble discret de $\C$, on peut choisir $\alpha$ de norme minimale. 
	
	On montre par l'absurde que $\LR = \left[ a,\alpha\right]$ : soit $z\in\LR$ qui n'est pas dans le réseau $\left[ a,\alpha\right]$. On peut choisir $z$ de norme minimal. Si $z$ est entier, $z\in \LR\cap\Z = a\Z$ ce qui est exclus. Donc $z$ est linéairement indépendant de $a$. Or $|z| \leq |\alpha| + |z - \alpha|$. Si $z\neq\alpha$, $|z| < |\alpha|$, ce qui contredit la minimalité de $|\alpha|$. 
	
	Donc $\LR = \left[ a, \alpha\right]$. Si l'on note $\alpha = b + cfw_{\K}$, $\alpha$ n'étant pas entier on a $c \neq 0$. De plus $\LR = \left[ a, ka + \alpha\right]$ pour tout $k\in\Z$, donc on peut choisir une base telle que $0 \leq b< a$. 
\end{proof}

\begin{Def}
	Soit $\K$ un corps quadratique imaginaire. Il admet un unique automorphisme de corps non trivial. Pour tout $\alpha\in\K$, on note $\bar{\alpha}$ son conjugué dans $\K$. Pour tout idéal $\LR$ d'un ordre $\OR$ de $\K$, on note $\bar{\LR} = \left\lbrace \bar{\alpha}, \alpha\in\LR\right\rbrace $
\end{Def}

\begin{The}\cite{Cox}[Prop 5.5,Cor 5.6]
	
	$\OR_{\K}$ est un anneaux de Dedekind, c'est à dire :
	\begin{enumerate}
		\item $\OR_{\K}$ est intégralement clos
		\item $\OR_{\K}$ est noethérien
		\item Ses idéaux premier sont maximaux 
	\end{enumerate}
	En particulier tout idéal $\LR$ admet une unique factorisation par des idéaux premier (à permutation prêt des facteurs)
	\begin{equation*}
		\LR = \PR_{1}^{e_{1}}\PR_{2}^{e_{2}}\ldots\PR_{g}^{e_{g}}
	\end{equation*}
	avec $e_{1}$, ... $e_{g}\in\N$ et $g\in\N$.
\end{The}

Un ordre $\OR\subset\K$ quelconque n'est pas un anneaux de Dedekind, mais seule la clôture intégrale pose problème.

\begin{Prop}\cite{Cox}[5.A, Ex 7.4]
	
	Un ordre $\OR$ d'un corps de nombre $\K$ est noethérien, et ses idéaux premiers sont maximaux.
\end{Prop}


Dans la suite, nous allons préciser l'étude de $\OR_{\K}$, généraliser certains résultats obtenues aux ordres quelconques, puis déduire un résultat de factorisation d'idéaux d'ordre de corps quadratiques imaginaires.





\subsubsection{Idéaux fractionnaires et groupes de classes d'idéaux}


Les ordres de corps quadratique imaginaire sont des anneaux "presque" principaux. Pour définir proprement cela, on introduit une relation d'équivalence sur les idéaux. Deux idéaux sont équivalent si et seulement s'ils sont égaux à un produit prêt par un idéal principal. Cependant énoncé ainsi la relation ne sera pas symétrique. On est alors amené à définir un notion plus large d'idéaux, les idéaux fractionnaire.

\begin{Def}
	Soit $\K$ un corps quadratique imaginaire et $\OR\subset\K$ un ordre. On définit une relation d'équivalence entre idéaux $\LR,\;\LR'\subset\OR$.
	\begin{equation*}
		\LR\simeq\LR'\iff\LR = \alpha\LR', \alpha\in\K.
	\end{equation*} 
	Un idéal fractionnaire de $\OR$ est un $\OR$-module de $\K$ de type fini, non nul. Les idéaux de $\OR$ seront appelé "idéaux entiers" pour les distinguer des idéaux fractionnaires. 
	On note $I(\OR_{\K})$ l'ensemble des idéaux fractionnaires de $\OR_{\K}$.
\end{Def}

\begin{Prop}\cite{Cox}[5.A, Ex 5.2]
	
	Les idéaux fractionnaires sont les ensembles de la forme $\alpha\LR$, où $\alpha\in\K$ et $\LR\subset\OR$ est un idéal entier de $\OR$.
\end{Prop}

Le nom d'idéal fractionnaire vient de l'existence d'une structure de groupe multiplicatif sur $I(\OR_{\K})$. On peut aussi généraliser les résultats de factorisation des idéaux entiers aux idéaux fractionnaires de $\OR_{\K}$

\begin{Prop}\cite{Cox}[Prop 5.7]
	 
	Les idéaux fractionnaire de $\OR_{\K}$ sont inversible, c'est à dire que pour tout $\IR\in I(\OR_{\K})$, il existe $\IR'\in I(\OR_{\K})$ tel que $(\IR)(\IR') = \OR_{\K}$.
	$I(\OR_{\K})$ est donc un groupe pour la loi multiplicative. 
	Tout idéal fractionnaire $\IR$ admet une unique factorisation par des idéaux entiers premiers ou leurs inverses (à permutation prêt des facteurs):
	\begin{equation*}
		\IR = \PR_{1}^{e_{1}}\PR_{2}^{e_{2}}\ldots\PR_{g}^{e_{g}}
	\end{equation*}
	avec $e_{1}$, ... $e_{g}\in\Z$ et $g\in\N$.
\end{Prop}

\begin{Prop}\cite{Cox}[Prop 5.7]\label{denomid}
	
	Soit $\IR$ un idéal fractionnaire d'un ordre quadratique imaginaire $\OR$. Il existe un entier $d$ tel que $\LR = d\IR$ soit un idéal entier. 
\end{Prop}


\begin{Def}
	Pour tout corps de nombre $\K$, $I(\OR_{\K})$ contient l'ensemble $P(\OR_{\K})$ des idéaux fractionnaires principaux, i.e de la forme $\alpha\OR_{\K}$ avec $\alpha\in\K$. 
\end{Def} 

\begin{Rem}
	L'inverse de $\alpha\OR_{\K}$ dans $I(\OR_{\K})$ est $\alpha^{-1}\OR_{\K}$ : $P(\OR_{\K}))$ est un sous-groupe de $I(\OR_{\K})$. Le quotient $I(\OR_{\K}/P(\OR_{\K}))$ pour la loi multiplicative correspond aux classes d'équivalences de la relation $\IR\simeq\IR'\iff\IR = \alpha\IR',\alpha\in\\K$.
\end{Rem}

\begin{Def}
	Pour tout corps de nombre $\K$, on appelle groupe de classes d'idéaux de $\OR_{\K}$ (où de $\K$) le quotient $Cl(\OR_{\K}) =  I(\OR_{\K})/P(\OR_{\K})$. Si le contexte le permet, on appellera $Cl(\OR_{\K})$ le groupe de classes de $\K$. 
\end{Def}


La notion de groupe de classes peut être défini pour tout anneaux de Dedekind, en considérant le corps de fraction. Une propriété remarquable de $\OR_{\K}$ comparé aux anneaux de Dedekind quelconques est que $Cl(\OR_{\K})$ est fini. 

\begin{Prop}\cite{Cox}[5.A]
	
	$Cl(\OR_{\K})$ est fini, on note $h(\OR_{\K})$ son cardinal. 
\end{Prop} 


On va maintenant définir une notion groupe de classes pour un ordre $\OR$ quelconque d'un corps quadratique imaginaire $\K$. La difficulté vient du fait que certains idéaux entiers d'un ordre $\OR$ ne seront pas inversibles. Si $\LR\subset\OR$ est un idéal, $\LR$ pourrait aussi être un idéal d'un ordre contenant $\OR$. On ne pourra définir une notion d'inverse pour $\LR$ qu'en le voyant comme idéal du plus gros anneaux possible, c'est à dire $\left\lbrace \alpha\in\K, \alpha\LR\subset\LR\right\rbrace$. 

\begin{Def}
	Soit $\OR$ un ordre d'un corps quadratique imaginaire $\K$. $\LR\subset\OR$ est un idéal propre de $\OR$ si c'est un idéal entier tel que
	\begin{equation*}
		\OR = \left\lbrace \alpha\in\K, \alpha\LR\subset\LR\right\rbrace
	\end{equation*} 
	$\IR$ est un idéal fractionnaire propre de $\OR$ s'il vérifie $\OR = \left\lbrace \alpha\in\K, \alpha\IR\subset\IR\right\rbrace$. On note $I(\OR)$ l'ensemble des idéaux fractionnaires propres.
\end{Def}

\begin{Prop}\cite{Cox}[7.4]
	
	Un idéal fractionnaire $\IR$ d'un ordre $\OR$ est inversible si et seulement s'il est propre.
\end{Prop}

\begin{Coro}
	Les idéaux fractionnaires principaux d'un ordre $\OR$ sont propres, et forment un sous groupe de $I(\OR)$.
\end{Coro}

\begin{Def}
	Soit $\OR$ un ordre quadratique imaginaire. On note $P(\OR)$ le sous groupe de $I(\OR)$ des idéaux fractionnaires principaux. Le groupe de classes d'idéaux de $\OR$ est le quotient $Cl(\OR) =  I(\OR)/P(\OR)$
\end{Def}

\begin{Prop}\cite{Cox}[7.A]
	
	$Cl(\OR)$ est fini, on note $h(\OR)$ son cardinal. 
\end{Prop} 

On sera amené à calculer dans le groupe $Cl(\OR)$. On peut déjà remarquer qu'il sera facile de trouver un idéal représentant l'inverse d'une classe d'idéaux grâce aux résultats suivant. On souhaite aussi avoir des résultats de factorisations d'idéaux entiers ou fractionnaires propres d'ordres quelconques.

\begin{Def}
	Soit $\K$ un corps quadratique imaginaire. Il admet un unique automorphisme de corps non trivial, appelé conjugaison. Pour tout $\alpha\in\K$, on note $\bar{\alpha}$ son conjugué dans $\K$. Pour tout idéal $\LR$ d'un ordre $\OR$ de $\K$, on note 
	\begin{equation}
		\bar{\LR} = \left\lbrace \bar{\alpha}, \alpha\in\LR\right\rbrace 
	\end{equation}
\end{Def}

\begin{Prop}\cite{Cox}[Prop 7.14]\label{NormesId}
	
	Soit $\OR$ un ordre d'un corps quadratique imaginaire $\K$, et $\LR\in I(\OR)$
	\begin{enumerate}
		\item $\forall\alpha\in\OR,\; N(\alpha\OR) = N_{\Q}^{\K}(\alpha)$
		\item $N(\LR\LR') = N(\LR)N(\LR')$ pour tout idéal entier propre $\LR'$.
		\item $\LR\bar{\LR} = N(\LR)\OR$, donc $\bar{\LR}$ est l'inverse de $\LR$ dans $Cl(\OR)$
	\end{enumerate}
\end{Prop}


Si $\LR\subset\OR$ est un idéal, le produit $\LR_{\K} = \LR\OR_{\K}$ fourni un idéal de $\OR_{\K}$. Inversement si $\LR_{\K}\subset\OR_{\K}$ est un idéal, l'intersection $\LR = \LR_{\K}\cap\OR$ est un idéal de $\OR$. On se restreignant à certains idéaux, ces deux applications seront réciproque l'une de l'autre et permettront de trouver des factorisations d'idéaux d'ordres quelconques. 

\begin{Def}\cite{Cox}[Lem 7.18]
	Soit $\LR\subset\OR$ un idéal entier d'un ordre quadratique imaginaire, et $f$ le conducteur de $\OR$. L'idéal $\LR$ est dit premier avec le conducteur si l'une des deux conditions équivalente suivante est vérifiée :
	\begin{enumerate}
		\item $N(\LR)$ est premier avec $f$
		\item $\LR + f\OR = \OR$
	\end{enumerate}
\end{Def}

\begin{Prop}\cite{Cox}[Lem 7.18]
	
	Si $\LR\subset\OR$ est un idéal premier avec le conducteur, alors il est propre. L'ensemble des idéaux premier avec le conducteur est stable par multiplication et forme un sous groupe de $I(\OR)$.  
\end{Prop}

\begin{Def}
	On note $I(\OR,f)$ le sous groupe de $I(\OR)$ engendré par les idéaux premier avec le conducteur, et $P(\OR,f)$ le sous groupe de $I(\OR,f)$ engendré par les idéaux principaux premier avec le conducteur.
\end{Def}

\begin{Prop}\cite{Cox}[Lem 7.19]
	
	Si $\OR$ est un ordre d'un corps quadratique imaginaire, de conducteur $f$, alors
	\begin{equation*}
		Cl(\OR) =  I(\OR, f)/P(\OR, f)
	\end{equation*} 
\end{Prop}

\begin{Rem}
	La proposition précédente implique en particulier qu'il existe dans toute classes d'idéaux un représentant premier avec le conducteur. Si $\OR = \OR_{\K}$, $f = 1$ et aucune restriction n'est faite. 
\end{Rem}

Dans la suite on fixe un ordre quadratique imaginaire $\OR$ de conducteur $f$, et on considère les idéaux de $\OR_{\K}$ de norme premier avec $f$.

\begin{Def}\cite{Cox}[7.C, Lem 7.18]
	Soit $f\in\N$. Un idéal entier $\LR_{\K}\subset\OR_{\K}$ est dit premier avec $f$ si l'une des deux conditions équivalentes suivantes est vérifiées :
	
	\begin{enumerate}
		\item $N(\LR_{\K})$ est premier avec $f$
		\item $\LR_{\K} + f\OR_{\K} = \OR_{\K}$
	\end{enumerate}
	
	On note $I_{\K}(f)$ le sous groupe de $I(\OR_{\K})$ engendré par les idéaux entiers premier avec $f$, et $P_{\K}(f)$ le sous groupe de $I_{\K}(f)$ engendré par les idéaux entiers principaux premier avec $f$
\end{Def}

\begin{The}\cite{Cox}[Prop 7.20]
	
	Soit $\OR$ un ordre d'un corps quadratique imaginaire $\K$, de conducteur $f$.
	\begin{enumerate}
		\item Si $\LR_{\K}\subset\OR_{\K}$ est un idéal premier avec $f$, alors $\LR_{\K}\cap\OR$ est un idéal de $\OR$ premier avec le conducteur, de même norme que $\LR_{\K}$.
		\item Si $\LR\subset\OR$ est un idéal premier avec le conducteur, alors $\LR\OR_{\K}$ est un idéal de $\OR_{\K}$ premier avec $f$, de même norme de que $\LR$.
		\item L'application $\LR\mapsto\LR\OR_{\K}$ induit un isomorphisme entre $I(\OR,f)$ et $I_{\K}(f)$ d'inverse $\LR_{\K}\mapsto\LR_{\K}\cap\OR$
	\end{enumerate}
\end{The}

Une première conséquence de ce théorème est l'application à la factorisation d'idéaux premiers avec le conducteur.

\begin{Coro}\cite{Cox}[7.C, Ex 7.26]
	
	Soit $\LR\subset\OR$ un idéal entier premier avec le conducteur $f$, alors $\OR/\LR \simeq \OR_{\K}/\LR\OR_{\K}$. En particulier l'application $\LR\mapsto\LR_{\K}$ conserve les idéaux premiers. 
\end{Coro}

\begin{Coro}\cite{Cox}[7.C, Ex 7.26]
	Un idéal $\LR\subset\OR$ premier avec le conducteur admet une unique factorisation (à permutation prêt des facteurs) par des idéaux premier, et premiers avec le conducteur.
	
	\begin{equation*}
		\LR = \PR_{1}^{e_{1}}\PR_{2}^{e_{2}}\ldots\PR_{g}^{e_{g}}
	\end{equation*}
	
	avec $e_{1}$, ... $e_{g}\in\N$ et $g\in\N$.
	
\end{Coro}

\begin{Coro}\cite{Cox}[7.C, Ex 7.26]
	
	Tout idéal fractionnaire $\IR \in I(\OR,f)$ admet une unique factorisation par des idéaux entiers premier appartenant à $I(\OR,f)$ (à permutation prêt des facteurs):
	\begin{equation*}
		\IR = \PR_{1}^{e_{1}}\PR_{2}^{e_{2}}\ldots\PR_{g}^{e_{g}}
	\end{equation*}
	
	avec $e_{1}$, ... $e_{g}\in\N$ et $g\in\N$.
	
\end{Coro}

Une autre conséquence du théorème précédent est la description du groupe de classes de $\OR$ à partir d'idéaux fractionnaires de $\OR_{\K}$ premier avec $f$.

\begin{Prop}\cite{Cox}[Prop 7.22]
	Si $\OR$ est un ordre d'un corps quadratique imaginaire, de conducteur $f$, alors $Cl(\OR) =  I_{\K}(f)/P_{\K}(f)$
\end{Prop}



\subsubsection{Décomposition des nombres premiers}


Dans $\Z$, l'idéal engendré par un nombre premier $p$ est un idéal premier. Cependant dans un ordre quadratique imaginaire $\OR$, l'idéal $p\OR$ peut ne pas être premier. On peut cependant le factoriser lorsque $p$ est premier avec le conducteur de $\OR$. 


\begin{Def}
	Soit $\OR$ un ordre d'un corps quadratique imaginaire $\K$. Soit $p$ un nombre premier, premier avec le conducteur. L'idéal $p\OR$ se factorise de façon unique comme produit d'idéaux premier avec le conducteur. Si un idéal premier $\PR$ apparaît dans cette factorisation, on dira que $\PR$ est au dessus de $p$. 
\end{Def}

Il est possible d'expliciter les idéaux $\PR_{i}$, d'après un théorème dû à Kummer. On énonce ce théorème dans le cas des corps quadratique imaginaire. Pour un énoncé général pour tout corps de nombre, on se réfère à \cite{Coh00}[Th 4.8.13].

\begin{The}{Théorème de Kummer, cas quadratique imaginaire}\label{Kummer}
	
	Soit $\K$ un corps quadratique imaginaire, et $\OR$ un ordre de $\K$ de discriminant $\Delta$. Soit $T(X) = X^2 - \Delta \in\F_{p}\left[ X\right] $. Soit $p$ un nombre premier, premier avec le conducteur de $\OR$
	\begin{enumerate}
		\item Si $T(X) = X^2$, alors 
		\begin{equation*}
			p\OR = \PR^{2}\; avec \; \PR = p\OR + \sqrt{\Delta}\OR.
		\end{equation*}
		
		\item Si $T(X)$ est irréductible, alors $p\OR$ est un idéal premier de norme $p^{2}$.
		
		\item Si $T(X) = T_{1}(X)T_{2}(X)$ se factorise de façon non trivial, alors 
		\begin{equation*}
			p\OR = \PR_{1}\PR_{2}\; avec \; \PR_{i} = p\OR + T_{i}(\sqrt{\Delta})\OR.
		\end{equation*}
		
	\end{enumerate}
\end{The}

\begin{Def}
	Soit $\OR_\Delta$ un ordre quadratique imaginaire. Soit un nombre premier $p$, premier avec le conducteur. Selon les trois cas possible du théorème précédent, on dira que $p$ est 
	\begin{enumerate}
		\item ramifié dans $\OR_\Delta$, si $\Delta = 0$ modulo $p$.
		\item inerte dans $\OR_\Delta$, si $\Delta$ n'est pas un carré modulo $p$.
		\item décomposé dans $\OR_\Delta$, si $\Delta$ est un carré modulo $p$.
	\end{enumerate}
\end{Def}

Ainsi connaître la factorisation de $p\OR_\Delta$ dans un ordre quadratique imaginaire $\OR_\Delta$ revient à calculer le symbole de Legendre $ \left( \frac{\Delta}{p}\right) $. Expliciter la factorisation dans le cas décomposé demande en plus de trouver les racines de $\Delta$ dans $\F_{p}$. 


\subsubsection{Corps de classe de Hilbert}

Soit $\K$ un corps de nombres (dans notre cas un corps quadratique imaginaire), et $\Lc/\K$ une extension galoisienne. On peut de même s'intéresser aux factorisations possibles des idéaux premier de $\OR_{\K}$ vu dans l'anneaux des entiers algébriques $\OR_{\Lc}$ (qui est encore un anneaux de Dedekind). 

\begin{Def}
	Soit $\K$ un corps de nombres, et $\PR_{\K}\subset\OR_{\K}$ un idéal premier. Soit $\Lc/\K$ une extension galoisienne. L'idéal $\PR_{\K}\OR_{\Lc}$ se factorise dans $\OR_{\Lc}$ :
	\begin{equation*}
		\PR_{\K}\OR_{\Lc} = \PR_{1}^{e_{1}}\PR_{2}^{e_{2}}\ldots\PR_{g}^{e_{g}}
	\end{equation*} 
	On note $g$ le nombre d'idéaux premiers de la factorisation. Les exposants $e_{i}$ sont appelés les indices de ramification de $\PR_{\K}$ dans $\Lc$.
	Puisque $\PR_{\K}\OR_{\Lc}\subset\PR_{i}$ et que $\PR_{i}$ est maximal, le corps $\OR_{\Lc}/\PR_{i}$ est une extension du corps $\OR_{\K}/\PR_{\K}$ et on note $f_{i}$ son degré. Les degrés $f_{i}$ sont appelés les degrés d'inerties de $\PR_{\K}$ dans $\Lc$.
\end{Def}


\begin{Prop}\cite{Cox}[Th 5.8, 5.9]
	
	Soit $\K$ un corps de nombres, $\Lc$ une extension galoisienne de $\K$ et $n = \left[ \Lc : \K\right] $. Soit $\PR_{\K}\subset\OR_{\K}$ un idéal premier. Alors les coefficients $e_{i}$ et $f_{i}$ de $\PR_{\K}$ dans $\Lc$ ne dépendent pas de l'indice $i$ i.e. de l'idéal premier $\PR_{i}\subset\OR_{\Lc}$ correspondant.  
	On note $e$ et $f$ leur valeurs. Alors $efg = n$.
\end{Prop}

\begin{Def}
	Soit $\K$ un corps de nombres, $\Lc$ une extension galoisienne de $\K$. On note $n = \left[ \Lc : \K\right] $. Un idéal premier $\PR_{\K}\subset\OR_{\K}$ est dit :
	\begin{enumerate}
		\item ramifié dans $\Lc$, si son indice de ramification $e$ vérifie $e > 1$.
		\item inerte dans $\Lc$, si son degré d'inertie $f$ vérifie $f = n$, ce qui implique $e = g = 1$ et $p\OR_{\K}$ maximal.
		\item décomposé dans $\Lc$, si $g = n$, ce qui implique $e = f = 1$.
	\end{enumerate}
\end{Def}

Le but de ce qui suit est d'introduire une extension particulière de $\K$, appelée le corps de classe de Hilbert de $\K$, dont le groupe de Galois est isomorphes au groupe de classes de $\OR_{\K}$.

\begin{Def} Soit $\Lc/\K$ une extension galoisienne d'un corps de nombre $\K$. Une injection $\sigma :\K\rightarrow\C$ est dite réelle si elle plonge $\K$ dans $\R$. Une injection réelle $\sigma$ de $\K$ est dite ramifiée dans $\Lc$ si elle peut se prolonger en une injection $\sigma :\Lc\rightarrow\C$ qui n'est pas réelle. 
\end{Def}

\begin{Rem}
	Les injections $\sigma :\K\rightarrow\C$ sont appelées les "premiers infinis" de $\K$. (voir \cite{Cox}[5.C])
\end{Rem}

\begin{Def}
	Une extension $\Lc/\K$ galoisienne d'un corps de nombre $\K$ est dit abélienne si son groupe de Galois est abéliens. Une telle extension est non-ramifiée si ni aucun idéal premier de $\OR_{\K}$ ni aucune injection réelle de $\K$ n'est ramifié dans $\Lc$.
\end{Def}

\begin{The}\cite{Cox}[Th 5.18]
	
	Soit $\K$ un corps de nombre. Il existe une unique extension $\HR/\K$ abélienne, non-ramifiée et maximale au sens de l'inclusion, c'est à dire telle que toute extension abélienne non-ramifiée de $\K$ est incluse dans $\HR$
\end{The}

Cette extension de $\K$ est liée au groupe de classe de $\OR_{\K}$ par le théorème de réciprocité d'Artin. L'énoncé exacte de ce théorème demande la construction du symbole d'Artin, qui généralise le symbole de Legendre pour les idéaux premiers (voir \cite{Cox}[Lem 5.19, Cor 5.21]). On énonce ici uniquement l'isomorphisme qui découle de ce théorème, et qui nous sera utile par la suite. 

\begin{The}{Théorème de réciprocité d'Artin pour le corps de classes}\cite{Cox}[Th 5.23]\label{CHilbert}
	
	Soit $\K$ un corps de nombre, $\HR$ son extension abélienne non-ramifiée maximale. Alors le groupe de Galois de l'extension $\HR/\K$ est isomorphe au groupe de classe $Cl(\OR_{\K})$
	\begin{equation*}
		Gal(\HR/\K)\simeq Cl(\OR_{\K})
	\end{equation*}
\end{The} 

\begin{Def}
	Soit $\K$ un corps de nombres. Son extension abélienne non-ramifiée maximale $\HR$ est appelé corps de classe de Hilbert de $\K$.
\end{Def}





\subsubsection{Forme standard d'un idéal et formes quadratiques binaires}

On énonce une correspondance entre d'une part les idéaux d'un ordre d'un corps quadratique imaginaire $\OR_{\Delta}$, et d'autre part certaines formes quadratiques définies positives de discriminant $\Delta$. Cela sera en particulier utile pour mener des calculs dans le groupe de classe. On commence par définir les formes quadratiques qui nous intéressent.

\begin{Def}
	
	Soit $f = ax^2 + bxy + cy^2$ une forme quadratique. $f$ est dite binaire si elle est définie positive et à coefficients dans $\Z$. Son discriminant est $\Delta = b^2 - 4ac$. On notera $f = f(a,b,c)$ pour désigner une forme quadratique binaire, ou $f = f(a,b)$ si le discriminant $\Delta$ est connu. $f$ est dite primitive si ses coefficients $a$, $b$ et $c$ sont premiers entre eux.
\end{Def}

On peut associer à une forme quadratique binaire de discriminant $\Delta$ un idéal de $\OR_{\Delta}$.

\begin{Prop}\cite{Coh00}[Th 7.7]
	
	Soit $f = f(a,b)$ une forme quadratique binaire de discriminant $\Delta$. L'application
	\begin{equation*}
		I : f(a,b)\mapsto\left[ a, \frac{-b + \sqrt{\Delta}}{2}\right] 
	\end{equation*}
	associe à $f$ un idéal de $\OR_{\Delta}$.
\end{Prop}

Pour définir la réciproque de l'application $I$, il faut pouvoir retrouver les coefficients $a$ et $b$ à partir d'une écriture adéquate de l'idéal $\LR$.

\begin{Prop}\cite{Buchmann95}[2.5] \label{formestandard}
	
	Tout idéal entier $\LR\subset\OR_{\Delta}$ s'écrit 
	\begin{equation*}
		\LR = m\left[ a , \frac{-b + \sqrt{\Delta}}{2}\right],
	\end{equation*}
	avec $m,a \in\N$, $b\in\Z$, $b^{2} = \Delta [4a]$, $-a<b\leq a$ et $c = (b^2 -\Delta)/4a$ avec $a, b, c$ premiers entre eux.
\end{Prop}

\begin{Rem}
	La proposition citée \cite{Buchmann95}[2.5] est énoncée dans le cas des corps quadratique réel, c'est à dire de discriminant positif. Mais à part une condition différente sur le coefficient $b$, la preuve s'adapte sans difficulté au cas imaginaire.
\end{Rem}

\begin{Def}
	Un idéal $\LR$ est donné sous forme standard s'il est donné par l'unique base définie par \ref{formestandard}. Lorsque $\Delta$ est connue, on notera $\LR = id(m,a,b)$ une écriture standard de $\LR$, selon les notations de \ref{formestandard}. On dira que $\LR$ est primitif si $\LR = id(1,a,b)$.
\end{Def}

\begin{Rem}
	On peut définir une écriture standard d'un idéal fractionnaire $\IR$ de $\OR_{\Delta}$. En effet d'après la proposition \ref{denomid}, il existe $d\in\N$ minimal tel que $d\IR = \LR$ soit entier, de forme standard $\LR = id(m,a,b)$. Alors la forme standard de $\IR$ est
	\begin{equation*}
		\LR = \frac{m}{d}\left[ a , \frac{-b + \sqrt{\Delta}}{2}\right],
	\end{equation*}
	et sera notée $id\left( \frac{m}{d},a,b\right) .$
\end{Rem}


\begin{Coro}
	Soit $\LR = id(m,a,b)$ un idéal entier de $\OR_{\Delta}$. Soit $c = \frac{\Delta - b^2}{4a}$. L'application
	\begin{equation*}
		F : id(m,a,b)\mapsto f(a,b) = ax^2 + bxy + cy^2
	\end{equation*}
	associe à $\LR$ une forme quadratique binaire de même discriminant. 
\end{Coro}

\begin{Prop}\cite{Cox}[Th 7.7]
	Les applications $I$ et $F$ induisent une bijection entre les idéaux entiers primitifs et les formes quadratiques binaires primitives.
\end{Prop}

On peut retrouver facilement la norme et l'inverse d'un idéal entier donné sous forme standard.

\begin{Prop}\cite{Coh00}[5.2.5]
	
	Soit $\LR = id(m,a,b)\subset\OR_{\Delta}$ un idéal entier sous forme standard. Soit $\LR' = id\left( \frac{m}{N(\LR)},a,-b\right) $ (qui est un idéal fractionnaire de $\OR_{\Delta}$). Alors
	\begin{enumerate}
		\item $N(\LR) = m^2a$
		\item $\LR\LR' = \left[ \gcd(a,b,c),\;\frac{-b + \sqrt{\Delta}}{2}\right]  $
		\item $\LR$ est propre si et seulement s'il est primitif, et dans ce cas $\LR\LR' = \OR_{\Delta}.$
	\end{enumerate}
\end{Prop}

\begin{Def}
	Deux formes quadratiques binaires sont équivalentes si leurs images par $I$ sont dans la même classe. On note $Cl(\Delta)$ les classes d'équivalences de formes quadratiques binaires de discriminant $\Delta$.
\end{Def}

Ainsi la restriction de l'application $F$ à $I(\OR_{\Delta})$ induit une structure de groupe sur les formes quadratiques binaires primitives. Ce morphisme passe au quotient et induit un isomorphisme de groupe entre $Cl(\OR_{\Delta})$ et $Cl(\Delta)$. Cependant la notion d'équivalence entre formes quadratiques binaires peut être définie sans faire le lien avec les idéaux de $\OR_{\Delta}$, ainsi qu'une loi de groupe sur les formes et sur leurs classes, via le produit le Dirichlet \cite{Cox}[3.8].

\begin{Def}
	Soit $f = (a,b)$, $f' = (a', b')$ deux formes quadratiques binaires de même discriminant $\Delta$. On défini le produit de Dirichlet $f*f'$ de la façon suivante : Soit $B$ défini modulo $2aa'$ par les relations suivantes
	\begin{enumerate}
		\item $B = b [2a]$
		\item $B = b' [2a']$
		\item $B^2 = \Delta [4aa']$
	\end{enumerate}
	Alors $f*f' = (aa', B)$ est une forme quadratique binaire de discriminant $\Delta$.
\end{Def}

\begin{Prop}\cite{Cox}[Th 7.7]
	
	Avec les notation précédentes, $I(f*g) = I(f)I(g)$ et $F(\LR\LR') = F(\LR)*F(\LR')$. Les applications $I$ et $F$ induisent des isomorphismes réciproques entre $Cl(\OR_{\Delta})$ et $Cl(\Delta)$.
\end{Prop}

La proposition suivante donne plusieurs moyens de trouver des formes équivalentes à une forme quadratique binaire donnée. On peut alors chercher une forme équivalente aux coefficients les plus petits possibles. 

\begin{Prop}\cite{Cox}[Thm 2.8]\label{actionMatriceForme}
	Les changements de variables en $(x,y)$ suivant appliqués à une forme quadratique binaire $f(a,b,c)$ permettent d'obtenir $f' = f(a',b',c')$ équivalente à $f$. 
	\begin{enumerate}
		\item $(x,y)\mapsto (x - qy, y),\;\forall\;q\in\Z$
		\item $(x,y)\mapsto (x + y, y)$ 
		\item $(x,y)\mapsto(-y,-x)$  
	\end{enumerate}
	Ces changements de variables correspondent à l'action de certaines matrices de $SL_{2}(\Z)$ sur les formes quadratiques binaires (voir \cite{Cox}[Eq 2.2]). 
\end{Prop}

\begin{Def}
	Soit $f = f(a,b,c)$ une forme quadratique binaire de discriminant $\Delta < 0$. $f$ est une dite réduite si $|b| \leq a \leq c$ et si de plus lorsque $a = c$, $b \geq 0$. 
\end{Def}

\begin{Prop}\cite{Cox}[2.8]\label{uniciteReduite}
	
	Soit $\OR_{\Delta}$ un ordre d'un corps quadratique imaginaire. Toute classe de $Cl(\Delta)$ contient une unique forme quadratique réduite.  
\end{Prop}

\begin{Prop}\cite{Cox}[2.12]
	Une forme quadratique réduite $f = f(a,b)$ vérifie 
	\begin{equation*}
		|b|\leq|a|\leq\sqrt{\frac{|\Delta|}{3}}.
	\end{equation*}
	On peut donc représenter chaque classe $Cl(\OR_{\Delta})$ en $O(\log(|\Delta|))$ bits. 
\end{Prop}

La preuve de la proposition \ref{uniciteReduite} fournit un algorithme donnant l'élément réduit de la classe d'une forme quadratique binaire donnée, en $\OR(\log(\Delta)^{2})$ opération dans $\Z$. Les différentes étapes de l'algorithme utilise les changements de variables cités dans la proposition \ref{actionMatriceForme}. 

\begin{algorithm}[h]
	\caption{Réduction de forme quadratique binaire}
	\label{ALGO reduction}
	\begin{algorithmic}[1]
		\Require Une forme quadratique binaire $(a,b,c)$
		\Ensure Une forme quadratique réduite équivalente.
		\State $test \gets False$
		\While{test = False}
		\If{ $|b| > a$}
		\State Calculer $q$, $r$ tels que $b = 2aq + r$, $-a<r\leq a$
		\State $c \gets c - bq + aq^2$
		\State $b \gets r$
		\EndIf
		\State $test \gets True$
		\If{$a > c$}
		\State Echanger $a$ et $c$
		\State $b \gets -b$
		\State $test \gets False$
		\EndIf
		\If{$b < 0$}
		\If{$a = -b$ OU $a = c$}
		\State $b \gets -b$
		\EndIf
		\EndIf
		\EndWhile
		\State \Return $(a,b,c)$
	\end{algorithmic}
\end{algorithm}


\begin{Rem}
	On peut vérifier avec l'algorithme \ref{ALGO reduction} que deux idéaux entiers sont dans la même classe : il suffit de réduire leurs formes quadratiques binaires associées et vérifier que l'on trouve le même résultat. 
\end{Rem}


Il sera utile par la suite de trouver des générateurs d'idéaux (primitifs) principaux à partir de leur forme standard. On peut utiliser l'algorithme \ref{ALGO reduction} pour cela. Soit $id(1,a,b)$ un idéal supposé principal, et soit $\beta$ un générateur. D'après \ref{NormesId}, $N((1,a,b)) = a = N_{\Q}^{\K}(\beta)$, donc étant donné la forme standard, on connaît la norme d'un générateur. 

\begin{Def}
	Un entier $m\in\N$ est dit représenté par une forme quadratique binaire $f$ s'il existe $s$, $t$ deux entiers tels que $m = f(s,t)$.
\end{Def}

\begin{Prop}\cite{Cox}[Th 7.7]
	
	Un entier $m\in\N$ est représenté par une forme quadratique binaire primitive $f$ si et seulement s'il existe un idéal $\LR\sim I(f)$ tel que $N(\LR) = m$.
\end{Prop}

\begin{Prop}
	
	Soit $\LR = [a , \frac{-b + \sqrt{\Delta}}{2}]$ un idéal de $\OR_{\Delta}$ primitif sous forme standard. Soit $f = F(\LR)$ sa forme quadratique binaire primitive associée. Alors 
	\begin{equation*}
		f(x,y) = ax^2 + bxy + cy^2 = N_{\Q}^{\K}\left( ax - y\frac{-b + \sqrt{\Delta}}{2}\right) / N(\LR)
	\end{equation*}
\end{Prop}

\begin{proof}
	On sait que par définition de $N_{\Q}^{\K}$, 
	
	\begin{align*}
		N_{\Q}^{\K}\left( ax - y\frac{-b + \sqrt{\Delta}}{2}\right) / N(\LR) &= \left( ax + y\frac{b}{2} + y\frac{\sqrt{\Delta}}{2}\right) \left( ax + y\frac{b}{2} - y\frac{\sqrt{\Delta}}{2}\right)/ a\\
		&= \left( \left( ax + y\frac{b}{2} \right)^{2} - \left( y\frac{\sqrt{\Delta}}{2}\right)^{2}\right)/a\\
		&=  ax^2 + bxy + y^2\left( \frac{b^2 - \Delta }{4a}\right) \\
		&= ax^2 + bxy + cy^2 = f(x,y),
	\end{align*}
	
	la dernière étape étant dû au fait que $f$ est de discriminant $\Delta = b^2 - 4ac$.
	
\end{proof}

\begin{Coro}
	Soit $\LR = (1,a,b)\subset\OR_{\Delta}$ un idéal primitif sous forme standard. $\LR$ est principal si et seulement si $1$ est représenté par $f$, c'est à dire s'il existe deux entiers $x, y\in\Z$ tels que $f(x,y) = 1$. De plus dans ce cas,
	\begin{equation*}
		\LR = \left( ax - y\frac{-b + \sqrt{\Delta}}{2}\right) \OR_{\Delta}.
	\end{equation*}
\end{Coro}

\begin{proof}
	$\LR$ est un idéal principal si et seulement s'il existe un élément de $\LR$ de norme égale à $N(\LR)$. Puisque $\LR = [a , \frac{-b + \sqrt{\Delta}}{2}]$, tout élément $\beta$ de $\LR$ est de la forme
	\begin{equation*}
		\beta = \left( ax - y\frac{-b + \sqrt{\Delta}}{2}\right),\; x,y\in\Z.
	\end{equation*}
	Alors $N_{\Q}^{\K}(\beta) = N(\LR) \iff f(x,y) = 1$ d'après ce qui précède.
\end{proof}

On cherche donc $(x,y)\in\Z^2$ solution de $f(x,y) = 1$. La notion de forme réduite va nous permettre de conclure. 

\begin{Prop}\cite{Cox}[Thm 2.8, 2.9]
	
	Soit $r = ax^2 + bxy + cy^2$ une forme quadratique binaire primitive. Alors tout les entiers représentés par $r$ sont positifs. Si $r$ est réduite, alors $a$ est le plus petit entier représenté par $r$, avec $r(1,0) = a$.
\end{Prop}

Ainsi lorsque $f$ est une forme quadratique binaire qui représente $1$, sa forme réduite $r$ est de la forme $x^2 \pm xy + cy^2$. On en déduit l'algorithme \ref{ALGOidealprincipal} permettant de trouver le générateur $\beta$ d'un idéal principal. Ceci est l'algorithme utilisé par défaut dans Sage pour la recherche de générateur.  

\begin{algorithm}[h]
	\caption{Générateur d'un idéal principal primitif donné sous forme standard}
	\label{ALGOidealprincipal}
	\begin{algorithmic}[1]
		\Require $(1,a,b)$ un idéal primitif d'un ordre $\OR_{\Delta}$ supposé principal.
		\Ensure Un élément $\beta$ tel que $(1,a,b) = \beta\OR_{\Delta}$. 
		\State $f\gets ax^2 + bxy + cy^2$
		\State $r\gets$ résultat de l'algorithme \ref{ALGO reduction} appliqué à $f$
		\State $L\gets$ liste des changements de variables utilisés durant la réduction de $f$ (voir la proposition \ref{actionMatriceForme}).
		\State $(x,y)\gets(1,0)$
		\State $(x,y)\gets(x',y')$ obtenus en appliquant tout les changements de variables retenus dans $L$ dans le sens inverse d'apparition.
		\State \Return $(x,y)$
	\end{algorithmic}
\end{algorithm}

Les algorithmes \ref{ALGO reduction} et \ref{ALGOidealprincipal} sont polynomiaux en les coefficients et le discriminant de la forme donnée en entrée. Plus précisément on a le résultat suivant. 

\begin{Prop}\cite{Coh00}[5.4.3]\label{nbEtapesReduction}
	
	Soit $(a,b,c)$ une forme quadratique binaire primitive de discriminant $\Delta$. L'algorithme \ref{ALGO reduction} renvoie une forme de réduite de $(a,b,c)$ en utilisant au plus $2 + \log\left( \frac{a}{\sqrt{|\Delta|}} \right)$ divisions euclidiennes.
\end{Prop}

\begin{Coro}\label{COMPLEX reduction}
	Soit $(a,b,c)$ une forme quadratique binaire de discriminant $\Delta$. La complexité de l'algorithme \ref{ALGO reduction} appliqué à $(a,b,c)$ est en 
	\begin{equation*}
		O\left( \log(a)\log\left( \frac{a}{\sqrt{|\Delta|}} \right) +\log(b)\right).
	\end{equation*}
	En particulier si $b\leq 2a$ alors la complexité de l'algorithme \ref{ALGO reduction} est majorée en $O(log(a)^2)$.
	
\end{Coro}

\begin{proof}
	
	La complexité de chaque passage dans la boucle de la ligne $2$ est essentiellement dominée par l'étape de la ligne $4$, qui correspond à une division euclidienne de $b$ par $2a$ donc en $O(\log(b))$. Après le premier passage dans la boucle, $|b|<a$ et les futures divisions euclidiennes seront toutes en $O(\log(a))$. On conclut avec la proposition \ref{nbEtapesReduction}.
\end{proof}

\begin{Coro}\label{COMPLEX idealprincipal}
	Soit $(1,a,b) = \beta\OR_{\Delta}$ un idéal principal primitif de $\OR_{\Delta}$ donné sous forme standard (en particulier $a$ est la norme de $\beta$). L'algorithme \ref{ALGOidealprincipal} renvoie $\beta$ avec une complexité en
	\begin{equation*}
		O\left( \log(a)\log\left( \frac{a}{\sqrt{|\Delta|}} \right)\right).
	\end{equation*}
\end{Coro}


\begin{Rem}
	Certaines sources \cite{BrokerCharlesLauter}, \cite{JaoSoukharev} utilisent un algorithme dû à Cornacchia pour résoudre le problème de l'idéal principal. Ils énoncent une complexité en $O(\log(N(\beta)))$. Cependant aucune preuve n'est donné, ni aucune description exacte de l'algorithme utilisé. Après analyse du problème il semble que cette algorithme ne fonctionne pas dans tout les cas. On estime que l'algorithme \ref{ALGOidealprincipal} a une complexité satisfaisante. De plus par défaut Sage utilise aussi l'algorithme \ref{ALGOidealprincipal}. 
\end{Rem}


\subsection{Rappels sur les isogénies de courbes elliptiques}
%pas/peu de preuves
%structure volcan, normalisation, anneaux d'endomorphisme 

\subsection{Représentation d'isogénies par action du groupe de classes d'idéaux}
%preuves si elles sont pertinentes (remarques idéaux propres)

\subsection{Réduction de courbe elliptique et relèvement de Deuring}
%pas de preuve, cf Kohel. Tableaux ?

